% SPDX-License-Identifier: CC-BY-4.0
% Copyright (C) 2026 Fabian Januszewski, Abhijit (the Langlandscape team)
% Backup frame: the standing hypotheses of the defect formula, shortened from the report
% (mathematics/chapters/s00/setup.tex ss:s00-H = Hyp. hyp:s06-H; cases (F0), (F1), (F1') from s06/s06-m2.tex, Prop. s06-levM).
\begin{frame}{Standing hypotheses of the defect formula}
  \small
  Setting: $D=q$, $N=qM$; all hypotheses \stattag{hypotheses, assumed where stated}.
  \begin{itemize}\setlength{\itemsep}{2pt}
    \item \textbf{(H0)} $p\ge5$, $q\ne p$, $(M,pq)=1$, $p\nmid\varphi(N)$; in particular $q\not\equiv1\bmod p$.
    \item \textbf{(H1)} $\T$ is Hida's full ordinary cuspidal Hecke algebra of tame level $\Gamma_0(N)$ over $\Lambda$,
      localised at a \emph{non-Eisenstein} maximal ideal $\m$; $F$ a primitive ordinary $q$-new branch, $\lambda_F:\T\to I_F$,
      $I_F$ a normal local domain finite over $\Lambda$.
    \item \textbf{(H2)} $\bar\rho_\m$ is absolutely irreducible on $G_{\QQ(\sqrt{p^*})}$ and $p$-distinguished.
    \item \textbf{(P)} a Kim--Ota point of weight $k_P\in[3,p-1]$ with trivial character exists (or the weight-two variant).
    \item \textbf{(H0$^-$)} $q\not\equiv-1\bmod p$, or $\bar\rho_\m$ ramified at $q$.
    \item \textbf{(Lev)} alternatively to (H0$^-$), one of the level cases:
      \begin{itemize}\setlength{\itemsep}{0pt}
        \item (F0) $M=1$;
        \item (F1) $M>1$ squarefree, $\bar\rho$ unramified at every $\ell\mid M$ (so $N(\bar\rho)=1$), and $\ell\not\equiv-1\bmod p$ for every $\ell\mid M$;
        \item (F1$'$) as (F1), but with some $\ell\mid M$, $\ell\equiv-1\bmod p$ \stattag{in our reading of Kim--Ota}.
      \end{itemize}
  \end{itemize}
  {\footnotesize Outside (H0$^-$) and (Lev) the case $q\equiv-1\bmod p$ with $\bar\rho$ unramified at $q$ is \stattag{open}.}
\end{frame}
